\documentclass[UTF8]{ctexbeamer}

% 主题设置
\usetheme{Madrid} % 可以尝试其他主题如 Warsaw, CambridgeUS, default
\usecolortheme{default} % 可以尝试 seahorse, dolphin, beaver 等

% 如果需要使用一些额外的宏包可以放在这里
\usepackage{graphicx}
\usepackage{booktabs}
\usepackage{amsmath, amsfonts, amssymb}

% 包含必要的页码信息等
\setbeamertemplate{footline}[frame number]

% PPT 信息设置
\title{在这里输入您的毕业论文题目}
\subtitle{毕业设计（论文）答辩}
\author{答辩人：您的名字 \quad 指导教师：导师名字}
\institute{您的学院名称 / 大学名称}
\date{\today}

\begin{document}

% 封面页
\begin{frame}
    \titlepage
\end{frame}

% 目录页
\begin{frame}{目录}
    \tableofcontents
\end{frame}

%------------------------------------------------
\section{研究背景与意义}
%------------------------------------------------
\begin{frame}{研究背景与意义}
    \begin{itemize}
        \item \textbf{研究背景}：
        \begin{itemize}
            \item 此处简述课题的背景和当前面临的问题...
        \end{itemize}
        \vspace{0.5cm}
        \item \textbf{研究意义}：
        \begin{itemize}
            \item 简述本课题在实际应用中或理论上的意义...
        \end{itemize}
    \end{itemize}
\end{frame}

%------------------------------------------------
\section{国内外研究现状}
%------------------------------------------------
\begin{frame}{国内外研究现状}
    \begin{itemize}
        \item 相关工作 1
        \vspace{0.3cm}
        \item 相关工作 2
        \vspace{0.3cm}
        \item 现有工作的不足（从而引出本论文的研究动机）
    \end{itemize}
\end{frame}

%------------------------------------------------
\section{主要研究内容与核心方法}
%------------------------------------------------
\begin{frame}{主要研究内容与核心方法}
    这里可以插入您的总体架构图：
    
    %\begin{figure}[htbp]
    %    \centering
    %    \includegraphics[width=0.8\textwidth]{figures/architecture.png}
    %    \caption{系统总体架构 / 方法框架}
    %\end{figure}
    
    重点介绍论文所提出的创新点核心算法。
\end{frame}

%------------------------------------------------
\section{实验设计与结果分析}
%------------------------------------------------
\begin{frame}{实验设计与结果分析}
    \begin{block}{实验环境与数据集}
        此处可以简写实验配置。
    \end{block}
    
    \begin{block}{结果分析}
        可以插入表格或者数据图以直观展示论文的成果：
        \begin{table}
            \centering
            \begin{tabular}{lcc}
                \toprule
                方法 & 指标 A & 指标 B \\
                \midrule
                Baseline & 80.0\% & 85.0\% \\
                \textbf{本文方法} & \textbf{85.0\%} & \textbf{89.0\%} \\
                \bottomrule
            \end{tabular}
            \caption{实验结果对比}
        \end{table}
    \end{block}
\end{frame}

%------------------------------------------------
\section{总结与展望}
%------------------------------------------------
\begin{frame}{总结与展望}
    \begin{itemize}
        \item \textbf{论文总结}：
        \begin{itemize}
            \item 总结一：完成了哪些工作...
            \item 总结二：取得了哪些成果...
        \end{itemize}
        \vspace{0.5cm}
        \item \textbf{未来展望}：
        \begin{itemize}
            \item 尚存在的不足之处。
            \item 未来的研究优化方向。
        \end{itemize}
    \end{itemize}
\end{frame}

%------------------------------------------------
\begin{frame}
    \centering
    \Huge
    谢谢各位老师的聆听！\\
    \vspace{1cm}
    \Large Please Questions
\end{frame}

\end{document}
